\documentclass[article, backcover, french, nodocumentinfo]{upmethodology-document}
% Encoding
	\usepackage[utf8]{inputenc}
	%\usepackage[latin1]{inputenc}
	\usepackage[T1]{fontenc}

% Language
	%\usepackage[francais,english]{babel} % Second language = main language

% Show a summary of the layout of the current document with \layout.
	%\usepackage{layout}

% For easy management of document margins and the document page size.
	%\usepackage[top=2cm, bottom=1.8cm, left=1.8cm, right=1.8cm, head=14pt, foot=36pt]{geometry}

% Lets you change line spacing.
	%\usepackage{setspace}

% Euro symbol
	%\usepackage{eurosym}

% Fonts (include only one)
	%\usepackage{bookman}
	%\usepackage{charter}
	%\usepackage{newcent}
	\usepackage{lmodern}
	%\usepackage{mathpazo}
	%\usepackage{mathptmx}

% Enables typesetting of hyperlinks
	%\usepackage{url}
	\usepackage{hyperref}

% Verbatim environment
	%\usepackage{verbatim}
	%\usepackage{moreverb}
	\usepackage{fancyvrb}

% Code listing
	\usepackage{listings}

% To change header and footer of any page of the document.
	\usepackage{fancyhdr}

% Allows you to insert graphic files within a document.
	%\usepackage{graphicx}

% Allows figures or tables to have text wrapped around them.
	%\usepackage{wrapfig}

% Adds support for colored text.
	\usepackage{xcolor}

% Allows tables rows and columns to be colored, and even individual cells.
	%\usepackage{colortbl}

% Mathematics
	\usepackage{amsmath}
	\usepackage{amssymb}
	\usepackage{mathrsfs}
	%\usepackage{asmthm}
	%\usepackage{mathtools}
	%\usepackage{bm} % Greek letters in math mode

% Provides control over the layout of the three basic list environments: enumerate, itemize and description.
	%\usepackage{enumitem}

% Interface to sectioning commands for selection from various title styles
	%\usepackage[nobottomtitles]{titlesec}

% Highly customized stacking of objects, insets, baseline changes, etc.
	\usepackage{stackengine}

% Routines for constrained scaling and stretching of objects, relative to a reference object or in absolute terms
	\usepackage{scalerel}

% Provides control over the typography of the Table of Contents, List of Figures and List of Tables, and the ability to create new ‘List of ...’.
	%\usepackage{tocloft}
	%\usepackage{titletoc}

% Advanced bibliography handling.
	%\usepackage{bibtex}
	%\usepackage{biblatex}

% Allows customization of appearance and placement of captions for figures, tables, etc.
	%\usepackage{caption}

% Provides the multicols environment which typesets text into multiple columns.
	%\usepackage{multicol}

% This package simplifies the insertion of external multi-page PDF or PS documents.
	%\usepackage{pdfpages}

% Prints out all index entries in the left margin of the text.
	%\usepackage{showidx}

% It allows to define multiple floats (figures, tables) within one environment giving individual captions and labels in the form 1a, 1b.
	%\usepackage{subcaption}

% Lets you insert notes of stuff to do with the syntax \todo{Add details.}.
	%\usepackage{todonotes}

% Text Companion fonts, which provide many text symbols (such as baht, bullet, copyright, musicalnote, onequarter, section, and yen), in the TS1 encoding.
	\usepackage{textcomp}

% For floating elements placement
	\usepackage{float}

% Personnal packages
	\usepackage{packages/MagicListings}
	\usepackage{packages/AlgorithmDefinition}

% For more information about UPmethodology: https://www.ctan.org/pkg/upmethodology

%=======================================================================================================
%=============================================== Informations ==========================================
%=======================================================================================================

%%% Document Information and Declaration
\declaredocument{Jeu tactique: Pogo}{Projet d'IA41}{UTBM\_IA41\_Pogo\_P17}

%%% Abstract and Key-words
\setdocabstract[french]{Projet UTBM de l'UV IA41 du semestre de printemps 2017}
\setdockeywords[french]{UTBM, IA41, Pogo}

%%% Document Authors and Validators
\addauthorvalidator*[julien.barbier@utbm.fr]{Julien}{Barbier}{Étudiant en INFO02}
\addauthorvalidator*[kadir.ercin@utbm.fr]{Kadir}{Ercin}{Étudiant en INFO02}
\addauthorvalidator*[maxime.pinard@utbm.fr]{Maxime}{Pinard}{Étudiant en INFO02}

%%% Informed People
\addinformed*[fabrice.lauri@utbm.fr]{Fabrice}{Lauri}{Professeur de l'UV IA41}

%%% Copyright and Publication Information
\setcopyrighter{Julien Barbier, Kadir Ercin et Maxime Pinard}
\setpublisher{Julien Barbier, Kadir Ercin et Maxime Pinard}
\setprintingaddress{France}

%%% Version
\incversion{\makedate{\the\day}{\the\month}{\the\year}}{Initial version.}{\upmpublic}

%=======================================================================================================
%================================================== Configs ============================================
%=======================================================================================================

% Change Front Page Layout
%\setfrontcover{modern} % modern or classic

% Change Illustration Picture
%\setfrontillustration[1.3]{figures/figure}

% Source code formatting
\upmcodelang{cpp} % uml, java or cpp

% Prevent page breaks in paragraphs
\predisplaypenalty=1000
\postdisplaypenalty=1000
\clubpenalty=1000

% Minimal space required in the bottom margin not to move the title on the next page
%\renewcommand{\bottomtitlespace}{.1\textheight}

% Links config, especialy for the table of contents
\hypersetup{
    colorlinks=true,
    linkcolor=black,
    urlcolor=blue,
    linktoc=all
}

% French language config
%\frenchbsetup{StandardLayout=true,ReduceListSpacing=false,CompactItemize=false}

%=======================================================================================================
%================================================= Functions ===========================================
%=======================================================================================================

%Paragraph with line break
\newcommand{\p}[1]{\paragraph{#1\\}}

% Function to print a warning sign
\newcommand{\dangersign}[1][2.5ex]
	{\renewcommand{\stacktype}{L}
		{\scaleto{\stackon[1pt]{\color{red}$\triangle$}{\fontsize{4pt}{4pt}\selectfont !}}{#1}}}

% Definition of some dt/dx/dy shortcuts for integrals
\newcommand{\dt}
{\;\mathrm{d}\,t}

\newcommand{\dx}
{\;\mathrm{d}\,x}

\newcommand{\dy}
{\;\mathrm{d}\,y}

% Definition of \Witem for 'itemize' environment with a warning sign
\newcommand{\Witem}
{\item[\dangersign{}]}

% Definition of a Max function shortcut
\newcommand{\Max}[2][ ]
{\underset{#1}{\text{Max}}\,#2}

% Definition of a Min function shortcut
\newcommand{\Min}[2][ ]
{\underset{#1}{\text{Min}}\,#2}


\begin{document}
	\upmdocumentsummary{}
	\upmdocumentauthors{}
	%\upmdocumentvalidators{}
	\upmdocumentinformedpeople{}
	\upmpublicationpage{}
	\thispagestyle{empty}
	\tableofcontents{}
	\lstlistoflistings{}
	%\listoffigures{}
	\newpage{}

	\section{Sujet, règles du jeu}
		\subsection{Présentation}
		\subsection{Règles}
		\subsection{Objectif}
	\section{Réalisation}
		\subsection{Représentation du jeu}
			\subsubsection{pions}
			\subsubsection{Piles de pions}
			\subsubsection{Plateau}
		\subsection{Structure de l'IA}
			\subsubsection{MinMax}
				\lstinputalgo[caption=MinMax]{algorithms/MinMax.algo}
			\subsubsection{AlphaBeta}
			\subsubsection{Génération des coups}
			\subsubsection{Fonction d'évaluation}
				\paragraph{}
					Afin d'évaluer le partie  
					nous comptons pour un joueur  le nombre de piles qu’il contrôle, que nous divisons par le nombre total de piles présentes sur le jeu.
					Notre évaluation représente donc le pourcentage de piles que nous contrôlons et Eval = 0 représente une défaite pour le joueur qui évalue car il ne possède aucune  et Eval = 1 une victoire car il les contrôle toutes. On a donc Eval $\in$ [0,1].
					\\\\De façon formelle :
					
					$$
					x_{i} = \left\{
					\begin{array}{ll}
						1 \mbox{ si  on contrôle la pile de la case i} \\
						0  \mbox{ sinon}
					\end{array}
					\right.
					$$
					
					et n le nombre de piles présentes sur le plateau
			    	$$Eval = \frac{1}{n}  \sum_{i = 1}^{9} x_{i} $$
				    \begin{center}
				    	 Exemple \\
				    \includegraphics{Eval.png}
					\end{center}
				\paragraph{}
					Dans ce cas pour le joueur violet contrôle
					4 piles, et le nombre de piles sur le plateau
					est de 7, ce qui nous donne Eval = $\frac{4}{7}$
					$\simeq$ 0,571

		\subsection{Interface utilisateur}
			\subsubsection{Réalisation}
			\subsubsection{Utilisation}
	\section{Résultats}
		\subsection{IA contre IA}
		\subsection{IA contre Humain}
%\begin{upmcaution}
%	This is an example of a caution message. This text must be rendered with enough height (usually 2 lines of text) to avoid intersection between %the caution icon and the box frame.
%\end{upmcaution}
%\begin{upminfo}
%	This is an example of an information message. This text must be rendered with enough height (usually 2 lines of text) to avoid intersection %between the caution icon and the box frame.
%\end{upminfo}
%\begin{upmquestion}
%	This is an example of a question message. This text must be rendered with enough height (usually 2 lines of text) to avoid intersection between %the caution icon and the box frame.
%\end{upmquestion}
%\begin{figure}[h!]
%	\centering
%	\includegraphics[width=0.5\textwidth]{figures/sample} % or <figures/sample.jpg>
%	\caption{sample figure}
%	\label{fig:sampleFigure}
%\end{figure}
\end{document}
